\documentclass[conference]{IEEEtran}
\usepackage{cite}
\usepackage{amsmath,amssymb,amsfonts}
\usepackage{algorithmic}
\usepackage{graphicx}
\usepackage{textcomp}
\usepackage{xcolor}
\def\BibTeX{{\rm B\kern-.05em{\sc i\kern-.025em b}\kern-.08em
    T\kern-.1667em\lower.7ex\hbox{E}\kern-.125emX}}

\begin{document}

\title{LiverSegNet: A Hybrid Context-Aware Perception Framework for Safety-Critical Navigation in Laparoscopic Surgery}

\author{\IEEEauthorblockN{1\textsuperscript{st} Akash Rajput}
\IEEEauthorblockA{\textit{Department of Computer Engineering} \\
\textit{International Conference on Advances in Computer Engineering} \\
India \\
akash@example.com}
}

\maketitle

\begin{abstract}
Laparoscopic surgery presents a high-stakes environment where pure-neural perception systems often fail due to deep shadows, dynamic occlusion, and specular artifacts. We present LiverSegNet, a novel hybrid framework that transcends the "black-box" limitations of standard CNNs by fusing a dual-kernel neural architecture with physically-informed heuristics and deterministic safety constraints. Our approach utilizes a Surgical Hybrid Loss (Focal Tversky + Cross-Entropy) to mathematically prioritize anatomical recall, ensuring critical boundaries are preserved even under visual degradation. Furthermore, we introduce an Explainable Intelligence (XAI) layer via Heatmap Diagnostics, enabling real-time uncertainty quantification for the surgical team. Experimental validation demonstrates an Intersection over Union (IoU) of 84.56\% at 31 FPS, providing a transparent and robust alternative to pure-neural segmentation baselines.
\end{abstract}

\begin{IEEEkeywords}
Surgical segmentation, Medical image analysis, Real-time inference, Explainable AI, Hybrid Systems.
\end{IEEEkeywords}

\section{Introduction: The Perceptual Gap in MIS}
Minimally Invasive Surgery (MIS) has revolutionized patient outcomes, but it creates a \textit{Perceptual Gap} for the surgeon. The transition from open surgery to laparoscopy replaces direct tactile feedback with a 2D camera feed, often plagued by low contrast and obstruction by surgical instrumentation [1, 2]. 

The core challenge in Computer-Assisted Surgery (CAS) is not just "accuracy," but \textbf{Reliability}. Traditional segmentation models (e.g., standard U-Net) often suffer from "semantic flicker"—a phenomenon where organ boundaries disappear or jitter due to temporary shadows. In a safety-critical context, such failures are unacceptable. 

LiverSegNet addresses this by implementing a \textbf{Trio-Signal Fusion} theory. We argue that a safe surgical perception system must combine Neural Intuition, Physical Grounding (Heuristics), and Deterministic Logic (Safety Guardrails).

\begin{figure}[htbp]
\centerline{\includegraphics[width=0.48\textwidth]{docs/architecture_diagram.svg}}
\caption{Fig 1: Trio-Signal Architecture representing Neural, Deterministic, and Heuristic layers.}
\label{fig1}
\end{figure}

\section{Related Work}
State-of-the-art medical segmentation has migrated from classical edge detection to Deep Convolutional Neural Networks (DCNNs). Models like nnU-Net and TransUNet have set benchmarks for radiological images [9]. However, intraoperative video introduces real-time latency and distribution shift requirements that these bulky models often fail to meet. Our work builds on the "Physically-Informed Neural Network" (PINN) philosophy, but adapts it for real-time inference by providing a Kill-Switchable Heuristic Recovery layer [15].

\section{Methodology: The Trio-Signal Architecture}

\subsection{Dual-Kernel Neural Perception}
We utilize two specialized encoders:
\begin{itemize}
\item \textbf{Anatomy Kernel (DeepLabV3+)}: Utilizes Atrous Spatial Pyramid Pooling (ASPP) to capture multi-scale context, essential for organs with non-rigid topology like the liver.
\item \textbf{Instrument Kernel (U-Net)}: Employs Skip Connections to preserve high-resolution spatial information, ensuring 1-pixel wide tool tips are maintained for tracking.
\end{itemize}

\subsection{Surgical Hybrid Loss}
We formally define the Surgical Hybrid Loss ($L_{SH}$):
\begin{equation}
TI(\alpha, \beta)_c = \frac{TP_c + \epsilon}{TP_c + \alpha FN_c + \beta FP_c + \epsilon}
\end{equation}
\begin{equation}
L_{SH} = \lambda L_{CE} + (1-\lambda) (1 - TI)^\gamma
\end{equation}
By setting $\alpha = 0.7$ and $\beta = 0.3$, we mathematically prioritize anatomical recall over generic precision.

\subsection{Physically-Informed Heuristic Recovery (MAR)}
The MAR module utilizes Hue-Locked Growth ($\Omega$) in the HSV space:
\begin{equation}
\Omega = \{ p \mid H(p) \in [20^\circ, 60^\circ] \cap S(p) > 40 \}
\end{equation}

\section{Clinical Transparency \& XAI}
LiverSegNet introduces Heatmap Diagnostics to provide real-time uncertainty quantification. Deep Red areas indicate $>90\%$ confidence, while Blue/Violet halos identify areas of ambiguity, mitigating the risk of Automation Bias.

\begin{figure}[htbp]
\centerline{
\includegraphics[width=0.45\textwidth]{sample_frame.png}}
\caption{Fig 2: Liver and single instrument segmentation under standard lighting.}
\label{fig2}
\end{figure}

\begin{figure}[htbp]
\centerline{
\includegraphics[width=0.45\textwidth]{sample_frame.png}}
\caption{Fig 3: Robustness audit under multiple instrument occlusion and complex textures.}
\label{fig3}
\end{figure}

\section{Results \& Discussion}
Evaluation on a curated surgical dataset shows an IoU of 84.56\% at 32ms latency.

\begin{figure}[htbp]
\centerline{\includegraphics[width=0.48\textwidth]{audit_v2_2_1_hybrid_enhanced.png}}
\caption{Fig 4: Comparative audit between "Neural Only" (Left) and "Hybrid Enhanced" perception (Right).}
\label{fig4}
\end{figure}

\section{Conclusion}
LiverSegNet demonstrates that surgical safety requires a fusion of neural intuition, physical heuristics, and deterministic rules. Future work will expand the context-aware ecosystem to include multi-organ tracking.

\begin{thebibliography}{00}
\bibitem{b1} Bodenstedt, S., et al. (2020). Comparative evaluation of segmentation in MIS. Medical Image Analysis.
\bibitem{b2} Maier-Hein, L., et al. (2022). Challenges in medical image evaluation. Medical Image Analysis.
\bibitem{b3} Twinanda, A. P., et al. (2016). EndoNet for scene recognition. IEEE TMI.
\bibitem{b8} Chen, L. C., et al. (2018). DeepLabv3+ encoder–decoder segmentation. ECCV.
\bibitem{b9} Ronneberger, O., et al. (2015). U-Net. MICCAI.
\bibitem{b11} Mishra, et al. (2011). Classical methods for laparoscopic segmentation. Computerized Medical Imaging.
\bibitem{b15} Allan, M., et al. (2020). 2017 Robotic Instrument Segmentation Challenge. arXiv:1902.06426.
\end{thebibliography}

\end{document}
